% ===================================================================
% CHAPTER 2: VERIFICATION TOOLS
% ===================================================================
\chapter{Verification Tools}
\label{ch:tools}

This chapter presents the hardware and simulation tools used to verify
the estimation and control results of the rest of this dissertation. It
separates what these tools inherit from prior lab work from what this
dissertation adds to them; Section~\ref{sec:ch2:summary} states that
division plainly.

\section{Decabots (Prior Work)}
% Ported from: thesis.tex Ch.2 Introduction (lines 593-622); Systems
% paper Sec. V-A hardware description; bib_sources.tex sys:52, sys:47.
\label{sec:ch2:decabots}

Every hardware result in this dissertation runs on the Decabot testbed,
a platform built by prior members of the Santa Clara University Robotic
Systems Laboratory and reused here without modification to the robots
or the motion capture system. The Decabot is a low-cost omnidirectional
rover developed for hardware verification of cluster-space control
schemes \cite{41}. A prior generation of this testbed encoded
vector field readings directly in RGB color and printed the result on
the floor as a map for the robots to sense, following an indoor testbed
feasibility study for vector-field adaptive navigation \cite{42}.
That RGB approach suffered a persistent limitation: interference between
color channels in the robots' undercarriage-mounted sensors degraded the
system's ability to recover accurate direction and magnitude
information, and the discrete color palette limited the effective
sampling rate to roughly 0.15~Hz, too slow to support the continuous,
closed-loop control laws developed in
Chapters~\ref{ch:first_order} and~\ref{ch:second_order}.

The remainder of this chapter describes what this dissertation adds on
top of this inherited platform: an HSV-based field encoding, a
neural-network sensor calibration that corrects for the channel
interference the RGB approach could not, a characterization of robot
dynamics and velocity response, and the simulation environment used
throughout the rest of this dissertation.

\section{Decabot Testbed and Motion Capture}
% Ported from: thesis.tex Ch.2 "Testbed Overview" (lines 662-671) and
% Control Architecture Overview (lines 645-661); Systems paper Sec. V-A
% for the workspace footprint used in current hardware trials; DETC2025-167604
% (thesis:33) Sec. 3.1 "Testbed Overview" and Sec. 5 "System Limitations"
% for the floor-map area and physical constraint figures.
\label{sec:ch2:testbed}

The physical testbed consists of three omnidirectional Decabot rovers
operating on 2.3~m\textsuperscript{2} multicolored printed floor maps,
an OptiTrack motion capture system providing centimeter-level position
accuracy via reflective markers on each robot's roof, and a central
computer running the control architecture in MATLAB/Simulink
\cite{39}. Robots communicate with the central computer over
WiFi at 10~Hz: at each control cycle the central computer receives
OptiTrack positions and sensor readings, computes the control laws, and
returns velocity commands, which each robot's onboard controller
converts into drive motor commands. Physical constraints include a
maximum 3~m~$\times$~6~m work area and approximately one hour of
battery life per robot \cite{39}; thesis.tex additionally reports
a lower bound of roughly 1~m~$\times$~1~m, a collision-clearance minimum
for a multirobot team. The three-robot hardware trials reported in
Chapter~\ref{ch:first_order} used a 1.6~m~$\times$~1.6~m printed map
within this range.

Each Decabot carries a downward-facing TCS34725 RGB color sensor,
enclosed in a light-shrouded housing to minimize external lighting
effects, whose channels are not fully isolated from one another per its
datasheet, the source of the interference problem
Section~\ref{sec:ch2:decabots} describes. Correcting for this
interference is the subject of Section~\ref{sec:ch2:nn_calibration}.

The robot controller layer, the cluster-space controller layer, and the
adaptive navigation layer that this testbed runs are described in full
in Chapter~\ref{ch:estimation}; this chapter is concerned only with the
physical and sensing tools those layers run on.

\section{Vector Field Encoding in HSV Color Space}
\label{sec:ch2:hsv}

\subsection{Limitations of RGB Encoding}
% Ported from: thesis.tex Ch.2 Introduction (lines 599-607); DETC2025-167604
% (thesis:33) Sec. 1 Introduction and Sec. 3.2 for detail thesis.tex dropped.

Encoding a vector field reading directly in RGB color, as the prior
testbed generation did, ties the sensor's three color channels directly
to the field's own two independent components (direction and
magnitude), with no channel held in reserve to detect or correct for
cross-talk between them. Approaches using red and green to represent
orthogonal vector components suffered from sensor channel interference
on the robots' RGB sensors, leaving only a small portion of the color
spectrum usable once high-interference regions were eliminated
\cite{39}. This interference made accurate decoding of direction
and magnitude information difficult, and was particularly severe for
vectors of small magnitude, where minor sensor value changes could
produce large apparent direction shifts \cite{39}. The TCS34725's
red, green, and blue channels are not fully isolated per its own
specification sheet \cite{39}, and this effect persisted even
after conversion to an HSV color space; traditional calibration methods
neither significantly improved performance nor supported effective
scaling across multiple robots \cite{39}, motivating the
supervised-learning calibration of Section~\ref{sec:ch2:nn_calibration}.
The resulting discrete color palette also limited sampling to roughly
0.15~Hz, too slow for continuous robot motion \cite{39}.

\subsection{Hue as Direction, Saturation as Magnitude}
% Ported from: thesis.tex Ch.2 "Modeling Vector Fields in Color Space"
% (lines 665-671); DETC2025-167604 (thesis:33) Sec. 3.2 for the
% value-versus-saturation reasoning thesis.tex dropped.
\label{sec:ch2:hue_sat}

A vector field reading has two components at each point: a direction
and a magnitude. This dissertation's testbed encodes direction in hue
and magnitude in saturation, allowing a two-dimensional vector field to
be represented as a printed image and sensed by a simple color sensor
rather than requiring active field generation. Hue runs from 0 to
$2\pi$ and represents flow direction; hue's circular topology is a
natural match for vector direction \cite{39}. Saturation, rather
than value, was chosen to encode magnitude, running from 0 to 1, after
early tests showed the value channel to be highly sensitive to printer
ink variation, ambient lighting, and shadow effects \cite{39}.
This separates the field's two physical components onto two different,
independently interpretable channels of the color space, rather than
overloading three RGB channels with two quantities.

\subsection{Printed Field Construction}
% Ported from: DETC2025-167604 (thesis:33) Sec. 3.2-3.3 for the field
% construction and scaling detail; Systems paper Sec. V-A (lines 328-345)
% for the current-work restatement of the same scaling.

This methodology was used to construct two vector fields for the
first-iteration testbed trials of \cite{39}: a fixed vortex and,
by adding a sinking component to the same vortex, a sinking vortex, both
centered on the printed map. Once computed on the print grid, vector
magnitudes were scaled to the 0.3--1.0 saturation range: values below
0.3 do not produce a commanded velocity strong enough to overcome
stiction (Section~\ref{sec:ch2:sat_velocity}), and values below 0.2 were
excluded from calibration training entirely
(Section~\ref{sec:ch2:nn_calibration}) \cite{39}. In the current
critical-points work, saturation is remapped to 0--1 in software, with
saturation 1.0 corresponding to the maximum robot velocity of 0.3~m/s
(the same cap given in Section~\ref{sec:ch2:dynamics}). Together with
the sensor shroud and the sensor's own illumination, this encoding is
minimally affected by overhead lighting variations; a test of turning
the overhead lab lights on and off confirmed no effect on readings from
the calibration palette \cite{39}. The robots' RGB sensors are
calibrated against a printed palette spanning hue 0 to $2\pi$ and
saturation 0 to 1 at 24 intervals in each dimension; each robot records
raw RGB and clear-channel counts over every palette cell. This printed-map
method has imperfections in both printing and sensing, addressed by the
calibration procedure of Section~\ref{sec:ch2:nn_calibration} and
revisited as a simulation noise source in
Section~\ref{sec:ch2:noise_model}.

\begin{figure}[h]
  \centering
  \includegraphics[width=0.75\textwidth]{ch2_measurement_error_comparison.png}
  \caption{Analytical fields, sensor-based reconstructions, and
  reconstruction error for vortex and saddle fields.}
  \label{fig:ch2:hsv_field}
\end{figure}
% FIGURE SOURCE: Paper_Writing/Systems Paper/cwaight_systems_paper/measurement_error_comparison.png

\section{Neural Network Sensor Calibration}
\label{sec:ch2:nn_calibration}
% Ported from: DETC2025-167604, "A Functional Indoor Testbed for
% Multirobot Adaptive Navigation in Vector Field Environments" (Waight
% and Kitts, 2025), cited throughout as thesis:33, Sec. 3.3-3.4 and
% Table 1, for the first-iteration methodology and numbers, verified
% directly against the extracted text at
% PhD Thesis/sources/idetc2025_testbed.txt; Systems paper Sec. V-A
% (lines 328-345) for the current, recalibrated numbers.
%
% Resolved: an earlier pass cited PhD Thesis/sources/idetc_paper.txt
% (DETC2024-142746, a control-primitives paper with no HSV/NN content)
% for this section by mistake; thesis:33 is the correct source and all
% numbers below are now checked directly against it.

\subsection{Calibration Palette and Data Collection}
% Ported from: thesis:33 Sec. 3.3.

The undercarriage color sensor's raw channels, red, green, blue, and a
clear channel, do not map cleanly to hue and saturation because the
channels are not fully isolated from one another. A calibration palette
with hues spanning 0 to $2\pi$ and saturations spanning 0 to 1, spaced at
24 equal intervals in each dimension, was printed, and three robots were
positioned above each palette cell to record RGB and clear-channel (K)
measurements alongside the intended hue and saturation values
\cite{39}. The three-robot dataset was combined and augmented by
adding $\pm 1\%$ noise to the RGB readings for model robustness. Raw
RGBK values were normalized to $[0,1]$, converted to hue, saturation,
and value using standard libraries, and appended to the original data,
giving a seven-dimensional input space $(R,G,B,K,H,S,V)$. Data points
with saturation below 0.2 were removed to improve model performance
(the same threshold noted in Section~\ref{sec:ch2:hsv}), and all inputs
and targets were normalized to $[-1,1]$ before training
\cite{39}.

\subsection{Network Architecture and Training}
% Ported from: thesis:33 Sec. 3.3.

This dissertation's solution is a pair of neural networks, trained per
robot (or, for a generalization test, trained on three robots and
validated on a fourth held out), that map the sensor's raw
seven-dimensional input $(R,G,B,K,H,S,V)$ to corrected hue and saturation
estimates. Both networks use two fully connected hidden layers with Tanh
activations; the saturation network has a single output neuron for
direct magnitude prediction, while the hue network has two output
neurons, sine and cosine of hue, to respect hue's circular topology
\cite{39}. A random search determined hidden layer sizes (2--12
units in the first layer, 4--10 in the second); patience and
augmentation parameters were also randomized, while the
Levenberg-Marquardt optimizer and default learning rate were held fixed.
Data was split 70/15/15 into training, validation, and test sets
\cite{39}. Candidate models were evaluated on eight validation
sets (the calibration and verification plots of all four robots,
Section~\ref{sec:ch2:generalization}); selection used a min-max
criterion, identifying each candidate's lowest $R^2$ score across all
eight sets and keeping the candidate with the highest such minimum, for
each network independently \cite{39}. The current calibration in
Chapter~\ref{ch:first_order}'s critical-points work uses a simplified
pair of feedforward networks, each with two hyperbolic-tangent hidden
layers and a four-dimensional $(R,G,B,K)$ input, with the saturation
network regressing a single output and the hue network outputting sine
and cosine of hue as above.

\subsection{Cross-Robot Generalization}
\label{sec:ch2:generalization}
% Ported from: thesis:33 Sec. 4.1 and Table 1, checked directly against
% the extracted paper text (PhD Thesis/sources/idetc2025_testbed.txt,
% lines 328-416). Replaces an earlier, coarser two-table summary drawn
% from thesis.tex that rounded off the per-dataset spread; the values
% below are the paper's own per-robot, per-dataset breakdown.

Summarizing across all validation sets, the cross-trained model achieves
an overall $R^2$ of 0.96 in direction (hue) and 0.91 in magnitude
(saturation) \cite{39}. The full per-robot, per-dataset breakdown
is given in Table~\ref{tab:ch2:calibration}: robots 1 through 3 supply
both a calibration (Cal) and verification (Ver) validation set each,
used in training the cross-robot model, while robot 4 (R4) is held out
of training entirely and used only to test generalization. Across the
six validation sets from robots 1--3, the cross-trained model's hue
$R^2$ ranges from 0.966 to 0.998 and saturation $R^2$ from 0.935 to
0.998. On the held-out robot 4, the cross-trained model achieves hue
$R^2 = 0.993$ (RMSE 0.024) and saturation $R^2 = 0.908$ (RMSE 0.073) on
its calibration set, and hue $R^2 = 0.964$ (RMSE 0.054) and saturation
$R^2 = 0.925$ (RMSE 0.079) on its verification set \cite{39}.
Individually tuned per-robot models improve further, typically to hue
$R^2 > 0.99$ and saturation $R^2 > 0.95$, at the cost of several hours
of per-robot calibration that the cross-trained model avoids
\cite{39}; individual models are not available for the held-out
robot 4 by construction.

\begin{table}[h]
\centering
\caption{First-iteration calibration model performance, cross-trained
versus individually tuned, on each robot's calibration (Cal) and
verification (Ver) validation set. Robot 4 was held out of training
entirely.}
\label{tab:ch2:calibration}
\begin{tabular}{llcccc}
\toprule
Robot-set & Model & Hue $R^2$ & Hue RMSE & Sat.\ $R^2$ & Sat.\ RMSE \\
\midrule
R1-Cal & Cross & 0.966 & 0.017 & 0.950 & 0.054 \\
       & Indiv & 0.999 & 0.009 & 0.972 & 0.040 \\
R1-Ver & Cross & 0.969 & 0.050 & 0.950 & 0.064 \\
       & Indiv & 0.995 & 0.020 & 0.963 & 0.055 \\
R2-Cal & Cross & 0.998 & 0.013 & 0.935 & 0.061 \\
       & Indiv & 0.996 & 0.018 & 0.968 & 0.043 \\
R2-Ver & Cross & 0.967 & 0.052 & 0.943 & 0.069 \\
       & Indiv & 0.991 & 0.027 & 0.952 & 0.063 \\
R3-Cal & Cross & 0.998 & 0.014 & 0.948 & 0.055 \\
       & Indiv & 0.999 & 0.007 & 0.980 & 0.034 \\
R3-Ver & Cross & 0.991 & 0.027 & 0.953 & 0.062 \\
       & Indiv & 0.997 & 0.015 & 0.976 & 0.044 \\
R4-Cal & Cross & 0.993 & 0.024 & 0.908 & 0.073 \\
       & Indiv & -- & -- & -- & -- \\
R4-Ver & Cross & 0.964 & 0.054 & 0.925 & 0.079 \\
       & Indiv & -- & -- & -- & -- \\
\bottomrule
\end{tabular}
\end{table}

\subsection{Recalibration After Maintenance}
% Ported from: Systems paper Sec. V-A (lines 328-345), the current,
% authoritative recalibration.

All robots were recalibrated after receiving preventive maintenance and
servicing for the critical-points work of Chapter~\ref{ch:first_order}.
This recalibration is the authoritative, current numbers referenced
throughout the rest of this dissertation; the numbers of
Table~\ref{tab:ch2:calibration} above are the first-iteration figures
they superseded. Using two feedforward networks trained on the
four-dimensional $(R,G,B,K)$ feature vector, with training pooled across
robots so a single model transfers to an uncalibrated unit without
per-robot tuning, the current calibration achieves hue RMSE of 0.17~rad
(9.7\textdegree, $R^2 = 0.991$) and saturation RMSE of 4.4\%
($R^2 = 0.951$) across the operational range
(Table~\ref{tab:ch2:recalibration}).

\begin{table}[h]
\centering
\caption{Sensor calibration accuracy, first iteration versus current
recalibration.}
\label{tab:ch2:recalibration}
\begin{tabular}{lcc}
\toprule
 & Hue & Saturation \\
\midrule
First iteration, cross-robot model overall & $R^2 = 0.96$ & $R^2 = 0.91$ \\
Current recalibration & RMSE 0.17 rad (9.7\textdegree), $R^2 = 0.991$ & RMSE 4.4\%, $R^2 = 0.951$ \\
\bottomrule
\end{tabular}
\end{table}

\begin{figure}[h]
  \centering
  \includegraphics[width=\textwidth]{ch2_color_sensor_hsv_predictions.png}
  \caption{Neural network predictions for HSV field encoding on
  unit-normalized axes. Hue RMSE = 0.027 (0.17~rad, 9.7\textdegree),
  saturation RMSE = 0.044 (4.4\%).}
  \label{fig:ch2:hue_predictions}
\end{figure}
% FIGURE SOURCE: Paper_Writing/Systems Paper/cwaight_systems_paper/color_sensor_hsv_predictions.png

\section{Saturation-to-Velocity Characterization}
% Ported from: thesis.tex Ch.2 "Robot Velocity from Saturation"
% (lines 729-736); DETC2025-167604 (thesis:33) Sec. 3.4 "Robot Velocity
% from Saturation Testing" for the methodology and the corrected
% threshold thesis.tex had rounded down to 0.2 (that is the separate
% training-data-exclusion threshold of Section~\ref{sec:ch2:hsv}, not
% the velocity-linearity threshold).
\label{sec:ch2:sat_velocity}

Because saturation encodes field magnitude, and a robot's commanded
velocity in the navigation primitives of Chapter~\ref{ch:estimation}
scales with sensed magnitude, this testbed also validates that robot
velocity responds linearly to sensed saturation above the stiction
floor (Section~\ref{sec:ch2:dynamics}). To test this, lanes of constant
hue at varying saturation were printed, each robot was placed on each
lane and commanded to follow the hue direction at the velocity indicated
by that lane's saturation, and each robot's position was tracked over a
3-second period per lane \cite{39}. Across all four robots,
velocity scaled linearly with saturation above 0.3; saturation values
below 0.3 did not produce a commanded velocity strong enough to overcome
static friction \cite{39}. This 0.3 velocity-linearity threshold
is separate from the 0.2 saturation threshold below which readings were
excluded from calibration training (Section~\ref{sec:ch2:hsv}).

\section{Robot Dynamics Identification}
\label{sec:ch2:dynamics}
% Ported from: Systems paper Sec. III-A (robot controller layer, lines
% 192-206) for the identified values; thesis.tex Ch.4 Sec.
% "Robot and Cluster Dynamics" (around lines 2015-2058) for the same
% values as previously drafted; Draft_11.tex Sec. III (lines 553-559)
% and Table I (lines 561-590) for the pentagon-code operating point;
% omnibot.py docstring and constructor defaults for the code values.
% The first-order lag MODEL equations belong to Chapter 4
% (\ref{ch:estimation}); this section reports only identification and
% measured values.

Each robot is modeled, in both simulation and the control law derivation
of Chapter~\ref{ch:estimation}, as a single omnidirectional point mass
with a first-order actuator lag response, discretized at control period
$\Delta t$ as
\begin{equation}
\mathbf{v}[k+1] = \alpha_{\text{mom}}\mathbf{v}[k] + (1-\alpha_{\text{mom}})\mathbf{v}_{\text{des}}[k],
\qquad \alpha_{\text{mom}} = e^{-\Delta t/\tau},
\label{eq:ch2:momentum}
\end{equation}
where $\tau$ is the actuator time constant identified on the Decabot
hardware. The full continuous-time model this discretization comes from
is given in Chapter~\ref{ch:estimation}.

Robots are limited to a maximum speed of 0.3~m/s, and require a minimum
commanded speed to overcome static friction (stiction) below which
commanded motion does not produce actual motion.

The robot parameters were identified twice, once for the testbed paper
\cite{39} and again after the sensors and drive were recalibrated for
the critical-points work of Chapter~\ref{ch:first_order}. The current
identification gives a time constant $\tau \approx 0.28$~s, so
$\alpha_{\text{mom}} \approx 0.70$ at the 10~Hz control rate; earlier
drafts used a rounder $\tau = 0.3$~s ($\alpha_{\text{mom}} \approx
0.717$). The simulator uses the fixed value $\alpha_{\text{mom}} = 0.7$.
%% Note: omnibot.py uses the constant 0.7, which does not rescale if dt
%% changes (its docstring audit note, 2026-07-03).

The stiction floor is $v_{\text{stiction}} = 0.025$~m/s after
recalibration, the value used by the simulator; the testbed paper's
characterization, before recalibration, reported a usable velocity range
of 0.05 to 0.3~m/s \cite{39}. Draft
11's pentagon-formation operating point (Table~\ref{tab:ch2:params})
carries the same 0.025~m/s stiction floor and 0.3~m/s speed cap for its
double-gyre benchmark, lowering stiction to 0.002 (in the field's own
non-dimensional length units) and setting $\alpha_{\text{mom}} = 0$ for
its ocean HFR trial, where one control cycle corresponds to ten
real-world minutes and any physically reasonable vessel time constant is
negligible by comparison.

\begin{table}[h]
\centering
\caption{Robot dynamics operating points identified on Decabot hardware
and used across simulation trials. Double-gyre and ocean columns are
the second-order operating points of Chapter~\ref{ch:second_order}, in
non-dimensional length/time units; the Decabot column is the identified
hardware value used by the first-order work.}
\label{tab:ch2:params}
\begin{tabular}{lccc}
\toprule
Parameter & Decabot (identified) & Double gyre & Ocean \\
\midrule
$v_{\max}$ & 0.3~m/s & 0.3 & 0.3 \\
$v_{\text{stiction}}$ & 0.025~m/s & 0.025 & 0.002 \\
$\alpha_{\text{mom}}$ & 0.70 & 0.7 & 0 \\
$\tau$ (s) & 0.28 & -- & -- \\
\bottomrule
\end{tabular}
\end{table}
%% Verified 2026-09-29 against Draft_11.tex Table I (tab:params). Double-gyre
%% and ocean columns are in non-dimensional length/time units.

\section{Field Reconstruction as a Simulation Noise Model}
\label{sec:ch2:noise_model}
% Ported from: Systems paper Sec. V-A (lines 344-345) and Sec. IV
% (lines 239-244) for the reconstructed-field noise source; thesis.tex
% Ch.5 "Noise Model" (around lines 4031-4057) for the two-channel
% measurement/position noise formulation; src/control/pentagon_primitives.py
% docstring (lines 393-411) for the current code-level noise hooks.

The printed-map encoding of Section~\ref{sec:ch2:hsv} has imperfections
in both printing and sensing. Rather than model this error abstractly,
the critical-points work of Chapter~\ref{ch:first_order} reconstructed
two fields (a vortex and a saddle) directly from robot-sensed
measurements collected on the hardware testbed, and used these
reconstructed fields, and the interpolation models fit to them, as a
realistic noise condition for simulation, alongside the six analytical,
noise-free canonical fields. This ties the simulator's "noisy" condition
to a measured quantity, the same sensed reconstruction error shown in
Fig.~\ref{fig:ch2:hsv_field}, rather than to an assumed noise
distribution.

The second-order (pentagon) simulation code generalizes this into two
independent, parametric noise channels rather than one measured
reconstruction. Measurement noise is additive on each robot's sensed
$(u,v)$ reading; position noise perturbs the location at which a robot
samples the field, leaving its true, controlled position untouched. In
the current simulator these are exposed as \texttt{cluster.measurement\_noise\_std}
and \texttt{cluster.position\_noise\_std} on the six-robot cluster
object, both defaulting to zero. The three-robot simulator has no
parametric channel; its noisy condition is the reconstructed field itself
(\texttt{NNField}, \texttt{RBFField}, and \texttt{BlendedField} in
\texttt{src/fields/field\_types.py}). The derivation of how each channel
propagates into the fitted coefficients is given in
Chapter~\ref{ch:second_order}.

%% Verified 2026-09-29: pentagon_primitives.py lines 417-418 read both
%% attributes with default 0.0; AnalyticalField takes no noise argument.

\section{Simulation Environment}
\label{sec:ch2:simulator}
% Ported from: thesis.tex Ch.5 "Simulation Program" (around lines
% 3934-3957); Systems paper Sec. IV (lines 239-244); src/robot/omnibot.py,
% src/simulation/runner.py, src/fields/field_types.py.

Simulation trials throughout this dissertation use a Python framework
under \texttt{trunk/Python\_Simulations/Vector\_Fields/VF\_Robot/}, whose
entry point is \texttt{experiments/main\_omni.py}, configured by editing
module-level constants rather than command-line flags. Each robot is
represented as an \texttt{Omnibot} object implementing the first-order
momentum dynamics of Eq.~\eqref{eq:ch2:momentum}
(Section~\ref{sec:ch2:dynamics}), together with the maximum-speed and
stiction constraints of Table~\ref{tab:ch2:params}. The integrator is
explicit Euler at the robots' native 10~Hz control rate
($\Delta t = 0.1$~s); a simulation loop
(\texttt{src/simulation/runner.py}) advances each robot's commanded
velocity, integrates position, and, for time-varying fields, advances
the field's internal clock once per control step.

Vector fields are represented through a common interface,
\texttt{get\_value(x, y)}, implemented by an \texttt{AnalyticalField}
wrapper around a closed-form field function for the canonical fields of
Appendix~\ref{app:fields}, and by neural-network and reconstructed-field
classes for the hardware-derived fields of
Section~\ref{sec:ch2:noise_model}. The same interface, together with the
noise channels of Section~\ref{sec:ch2:noise_model} and the cluster-space
kinematics of Chapter~\ref{ch:estimation}, is reused unchanged across the
three-robot and six-robot formations studied in this dissertation.

\section{Summary of Prior and New Contributions}
\label{sec:ch2:summary}
% New text, modeled on Hart (2023) Ch. 2.4's pattern of crediting prior
% tools before listing new contributions.

Table~\ref{tab:ch2:contributions} separates what this chapter's tools
inherit from prior Robotic Systems Laboratory work from what this
dissertation adds to them.

\begin{table}[h]
\centering
\caption{Prior versus new contributions among this dissertation's
verification tools.}
\label{tab:ch2:contributions}
\begin{tabular}{ll}
\toprule
Prior work & New in this dissertation \\
\midrule
Decabot rovers \cite{41} & HSV field encoding \\
Low-cost indoor testbed \cite{42} & Neural-network sensor calibration \\
OptiTrack motion capture & Post-maintenance recalibration \\
Three-layer control architecture & Saturation-to-velocity characterization \\
 & Robot dynamics identification ($\tau$, stiction) \\
 & Field-reconstruction simulation noise model \\
 & Simulation environment (Omnibot, noise channels) \\
\bottomrule
\end{tabular}
\end{table}

The platform, the Decabot rovers and the indoor testbed they operate in,
is prior work. This dissertation's contribution to that platform is the
HSV encoding that replaces direct RGB field representation, the
neural-network calibration that corrects for the resulting sensor's
channel interference, and the recalibration, dynamics identification, and
simulation tooling built on top of it. These tools are used without
further comment for the remainder of this dissertation: the estimation
framework of Chapter~\ref{ch:estimation} and the zeroth-, first-, and
second-order results of Chapters~\ref{ch:zeroth}
through~\ref{ch:second_order} all run on the testbed and simulator
described here.
